\section{Related Work}
\label{sec:related}

Prompt optimizers edit one text artifact. OPRO conditions a model on past (solution, score) pairs \citep{opro}, ProTeGi turns minibatch errors into textual gradients \citep{protegi}, and TextGrad passes textual feedback back through a compound system \citep{textgrad}. MIPROv2 tunes instructions and few-shot demonstrations for each module of a program \citep{miprov2-dspy}. GEPA reflects on rollout traces, keeps a per-task Pareto set of prompts, and merges candidates that share an ancestor \citep{gepa}. We run both from public code. CapEvolve edits a file set (system prompt, tool code, MCP tool descriptions, skill packages) rather than one string, so its search space is wider than any of these methods covers.

Skill and context methods store procedure in text the frozen agent reads. ACE grows a playbook through incremental delta updates to avoid rewriting it wholesale \citep{ace}. EvoSkill diagnoses failures, writes skill folders, and keeps a frontier of programs that score well on a validation set \citep{evoskill}. SkillOpt applies bounded edits to one skill document and accepts an edit only if it strictly improves a held-out selection score \citep{skillopt}. Trace2Skill consolidates many trajectories in parallel into one skill directory \citep{trace2skill}, and WikiSkill co-evolves skills with a persistent wiki \citep{wikiskill}. In the tables of GEPA and SkillOpt, each cell is one score without an interval. We run these methods under one optimizer model and one rollout budget, with three seeds for the main arms, and report intervals. WikiSkill has no official code, so our arm uses a community port.

Harness and agent search edits code. ADAS has a meta agent write new agents in code from an archive of earlier ones \citep{adas}. AFlow searches code-represented workflows with Monte Carlo tree search \citep{aflow}, and AlphaEvolve evolves programs against an automated evaluator \citep{alphaevolve}. The Darwin G\"odel Machine lets a coding agent edit its own code and screens children on 10, then 60, then 200 tasks \citep{dgm}; the Huxley-G\"odel Machine instead guides expansion by an estimate of how well an agent's descendants perform \citep{hgm}. Meta-Harness gives a coding-agent proposer a filesystem holding the code, scores and traces of every earlier candidate \citep{meta-harness}. HarnessLens verifies a candidate only on tasks relevant to its change \citep{xu2026harnesslens}, and DarwinX admits a variant only if it extends the solved set without regressions \citep{darwinx}. CapEvolve's candidate DAG, merge verification and agent-driven controller overlap with these designs, and we do not claim the ideas are new. We ask whether they help once every arm has the same optimizer, budget and final set.

Acceptance rules decide which edit survives. CAPO races candidates and drops weak ones early \citep{capo}, and TRIPLE casts prompt selection as fixed-budget best arm identification \citep{triple}. BudgetAPO sizes evaluation slices from measured task noise \citep{liu2026budgetapo}. SAGE compares candidate and incumbent on the same items with a one-sided paired test \citep{wang2026sage}. PACE uses an anytime-valid e-process, which controls false commits under repeated testing \citep{shawn2026pace}, and a further paper studies when a self-evolving system should stop \citep{yin2026enough}. These rules carry guarantees that our rule P lacks. Rule P (Section~\ref{sec:method:P}) is a paired screen followed by confirmation on fresh trials, tuned by simulation, and its confirmation cost is charged to the search budget.

Evaluation hygiene is an old problem. Reusing a holdout set adaptively makes it a training set \citep{adaptive-data-analysis-holdout-reuse}, and the Ladder limits what a leaderboard reveals for the same reason \citep{ladder}. For agent optimizers, Wang et al.\ report that harness evolution on Terminal-Bench 2.1 does not beat simple test-time scaling under matched feedback and inference budgets, and that earlier work searched on the tasks it reported \citep{rethinkharness}. Meta-Harness states that it searches and scores on the same 89 TerminalBench-2 tasks \citep{meta-harness}. A study of recursive self-evolution finds that removing held-out selection gives a perfect in-sample score and a 33-point drop on test \citep{nguyen2026rsea}. Single-lineage prompt optimization is reported to match or beat GEPA with fewer rollouts \citep{naiveprompt}, and one analysis pass of a coding agent over stored trajectories to beat SkillOpt on all four benchmarks tried \citep{codingagentsopt}. SkillsBench compares matched conditions with and without skills \citep{skillsbench}. We follow this line. Our design adds a scorer audit, gold files hidden from the optimizer, fresh final sets, matched budgets, and one-shot and seed controls.
